\chapter{Body Dysmorphic Disorder}

\section*{Definition and Overview}
Body dysmorphic disorder (BDD) is a mental health condition in which a person is preoccupied with a perceived flaw or defect in their appearance that others cannot see or consider trivial. The preoccupation causes significant distress and interferes with daily life, work, study, or relationships. Unlike ordinary dissatisfaction with one's looks, the concerns in BDD are repetitive, intrusive, and hard to resist, and they may dominate thinking for hours each day. The person often checks mirrors constantly, seeks reassurance, or tries to hide or fix the perceived flaw, and the condition responds well to psychological and medical treatment.

\section*{Epidemiology}
Body dysmorphic disorder is thought to affect roughly 1 in 50 people, with a similar prevalence in men and women. Onset is most often in adolescence, a life stage when appearance concerns are common and self-esteem is still developing. It is probably underdiagnosed, because people may be ashamed of their concerns and avoid seeking help, and they may present instead with depression, anxiety, or social withdrawal. People with BDD commonly also have another mental health condition, most often depression, social anxiety disorder, or obsessive-compulsive disorder.

\section*{Aetiology and Pathophysiology}
The exact cause of BDD is not fully understood, but it is likely to arise from a combination of factors:

\begin{itemize}
    \item \textbf{Genetics:} the disorder tends to run in families, and it is more common in people with a family history of BDD or obsessive-compulsive disorder
    \item \textbf{Brain chemistry and structure:} differences in the brain regions that process visual information, emotion, and decision-making, and in the serotonin system, appear to be involved
    \item \textbf{Psychological factors:} a tendency to focus attention on small details of appearance rather than the whole, and to interpret others' reactions as negative
    \item \textbf{Life experiences:} teasing, bullying, or criticism about appearance in childhood, and unrealistic ideals of beauty promoted by media and social networks
    \item \textbf{Personality traits:} perfectionism and low self-esteem are common among people with BDD
\end{itemize}

\section*{Clinical Features}
\begin{itemize}
    \item Preoccupation with one or more perceived flaws, most often involving the skin, hair, or nose
    \item Repetitive behaviours such as mirror checking, excessive grooming, or skin picking
    \item Constantly comparing one's own appearance with that of others
    \item Seeking reassurance about appearance from family or friends
    \item Attempts to hide the perceived flaw with clothing, make-up, or posture
    \item Avoiding social situations, work, or school because of appearance concerns
    \item Difficulty concentrating, low mood, or anxiety
    \item A strongly held belief that the flaw is real and obvious to others (in severe cases this can reach delusional intensity)
    \item Thoughts of self-harm or suicide in severe cases
\end{itemize}

\section*{Differential Diagnosis}
\begin{itemize}
    \item Normal dissatisfaction with appearance or fleeting appearance concerns
    \item Eating disorders such as anorexia nervosa or bulimia nervosa
    \item Social anxiety disorder, in which the fear is of being judged rather than of a specific flaw
    \item Obsessive-compulsive disorder, in which the obsessions are not appearance-related
    \item Depression with feelings of worthlessness that focus on appearance
    \item Anxiety about a real, visible condition such as severe acne, scarring, or a disfiguring illness
    \item Schizophrenia or other psychotic disorders when the belief is fixed and bizarre
\end{itemize}

\section*{When to Seek Medical Care}
See a GP or mental health professional if appearance concerns consume a significant amount of time, cause distress, or interfere with work, study, or social life. Because people with BDD often hide their concerns, it is important to seek help even if the worry feels embarrassing. If depression or anxiety are also present, both conditions should be treated together.

\textbf{Seek immediate help if there are thoughts of self-harm or suicide.} Call 000 (or local emergency services), contact Lifeline on 13 11 14, or go to the nearest hospital emergency department. Acute mental health crises are medical emergencies.

\section*{Diagnosis and Investigations}
\begin{itemize}
    \item \textbf{Clinical interview:} a thorough history of the appearance concerns, including how much time they occupy, the behaviours they cause, and their impact on daily life
    \item \textbf{Mental state examination:} to assess mood, anxiety, risk of self-harm, and whether the beliefs are fixed
    \item \textbf{Risk assessment:} direct questions about suicidal thoughts, plans, and past self-harm
    \item \textbf{Screening questionnaires:} tools such as the Body Dysmorphic Disorder Examination or a simple appearance-preoccupation questionnaire may help identify the disorder
    \item \textbf{Physical examination:} usually not needed, but may reassure that there is no medical cause for a real skin or other physical problem
    \item \textbf{Assessment of co-morbid conditions:} screening for depression, anxiety, obsessive-compulsive disorder, and eating disorders
\end{itemize}

\section*{Management}
\begin{itemize}
    \item \textbf{Cognitive behaviour therapy (CBT):} the main psychological treatment, which targets the unhelpful thoughts about appearance and the checking, avoidance, and reassurance-seeking behaviours that maintain the disorder
    \item \textbf{Exposure and response prevention:} a component of CBT in which the person gradually faces feared situations and stops compulsive behaviours such as mirror checking
    \item \textbf{Selective serotonin re-uptake inhibitors (SSRIs):} antidepressant medicines that reduce the intensity of the preoccupation; these require specialist prescribing and monitoring
    \item \textbf{Combined treatment:} for many people, CBT together with an SSRI works better than either alone
    \item \textbf{Family involvement:} education and support for family members, who may unknowingly reinforce the disorder through reassurance
    \item \textbf{Treatment of co-morbid conditions:} addressing any coexisting depression or anxiety
    \item \textbf{Follow-up:} regular review to monitor progress and prevent relapse; severe cases may need referral to a specialist psychiatric team
\end{itemize}

\section*{Complications}
\begin{itemize}
    \item Depression, which may develop as a result of the distress caused by the disorder
    \item Social isolation, relationship breakdown, and withdrawal from school or work
    \item Repeated and unnecessary cosmetic procedures, which rarely satisfy and may make the concern worse
    \item Skin damage or infection from repeated picking or grooming of the affected area
    \item Excessive use of make-up or clothing to hide the perceived flaw
    \item Self-harm and suicidal thoughts or behaviour -- the most serious complications
\end{itemize}

\section*{Prognosis and Outcomes}
Body dysmorphic disorder often persists for years if untreated, but many people improve substantially with appropriate treatment. CBT and SSRIs are effective for a large proportion of those affected, although symptoms may take weeks or months to respond. The disorder tends to run a chronic, fluctuating course, and some people need longer-term treatment or repeated courses of therapy. Early help-seeking is associated with a better outlook. Even with treatment, appearance concerns may not disappear entirely, but they can become manageable and stop dominating daily life.

\section*{Prevention and Self-Care}
\begin{itemize}
    \item Seek help early, before the preoccupation becomes well established
    \item Limit exposure to social media and advertising that promote unrealistic beauty standards
    \item Avoid mirror checking and reassurance-seeking; these behaviours tend to feed the preoccupation
    \item Develop interests and activities that build self-esteem independently of appearance
    \item Build a supportive social network and talk to trusted people about how you feel
    \item Challenge unhelpful thoughts by asking how you would view a friend with the same feature
    \item Avoid cosmetic procedures for a problem others cannot see; discuss any urge for surgery with a mental health professional first
    \item Follow the treatment plan and attend review appointments to reduce the chance of relapse
\end{itemize}

\section*{Sources and Further Reading}
The following sources provide reliable, up-to-date information on body dysmorphic disorder:

\begin{itemize}
    \item NHS. \textit{NHS guidance on body dysmorphic disorder.} https://www.nhs.uk/conditions/
    \item Mayo Clinic. \textit{Information on body dysmorphic disorder symptoms and treatment.} https://www.mayoclinic.org/
    \item MSD Manual. \textit{Overview of body dysmorphic disorder.} https://www.msdmanuals.com/
    \item MedlinePlus. \textit{Body dysmorphic disorder patient information.} https://medlineplus.gov/
    \item NICE. \textit{Obsessive-compulsive disorder and body dysmorphic disorder guidelines.} https://www.nice.org.uk/
    \item World Health Organization. \textit{Mental health and disorders information.} https://www.who.int/
\end{itemize}
